\section{Data and measurement}\label{sec:data}

\subsection{Sources}

All series are official, harmonised and quarterly unless stated otherwise. The core panel
covers the present-day \val{n_countries} member states, from 2008-Q1 (for lags and
differences) to 2026-Q2. It was downloaded from the Eurostat Statistics API on
\val{snapshot_date} and frozen; its manifest has SHA-256 \val{hash_quarterly}\dots.
Table~\ref{tab:data} lists the series; Eurostat datasets are cited by their codes.

\begin{table}[htbp]\centering\small
\caption{Data. SA: seasonally adjusted; SCA: seasonally and calendar adjusted; NSA: not
adjusted. All snapshots are frozen with SHA-256 hashes (Appendix~\ref{app:repro}).}
\label{tab:data}
\footnotesize
\begin{tabularx}{\textwidth}{>{\raggedright\arraybackslash}Xl
  >{\raggedright\arraybackslash}Xl}\toprule
Quantity (unit) & Dataset & Filters & Use \\ \midrule
Real GDP growth (\% q/q) & \citetalias{eurostat_namq_10_gdp} & B1GQ, CLV\_PCH\_PRE, SCA & all \\
Unemployment rate (\% of labour force) & \citetalias{eurostat_une_rt_q} & Y15-74, T, SA & all \\
Participation rate (\% of population) & \citetalias{eurostat_lfsi_emp_q} & ACT, Y15-74, SA & all \\
GDP deflator (\% q/q) & \citetalias{eurostat_namq_10_gdp} & PD\_PCH\_PRE\_NAC, SCA & SVAR \\
Unemployment, monthly (\%) & \citetalias{eurostat_une_rt_m} & TOTAL, SA & monthly LP \\
Industrial production (index) & \citetalias{eurostat_sts_inpr_m} & B-D, SCA, 2021=100 & monthly LP \\
Employment, domestic and resident (persons) & \citetalias{eurostat_namq_10_pe} & EMP\_DC, EMP\_NC, SCA
  & cases \\
Hours worked, total and finance & \citetalias{eurostat_namq_10_a10_e} & TOTAL, K, THS\_HW,
  SCA & cases \\
Value added, total and finance & \citetalias{eurostat_namq_10_a10} & B1G, TOTAL, K,
  CLV20\_MEUR, CP\_MEUR, SCA & cases \\
Temporary employees (persons) & \citetalias{eurostat_lfsq_etgaed} & Y15-74, NSA & cases \\
All employees (persons) & \citetalias{eurostat_lfsq_eegaed} & Y15-74, NSA & cases \\
Temporary and part-time shares (\%) & \citetalias{eurostat_lfsq_etpga}, \citetalias{eurostat_lfsq_eppga} & Y15-74,
  NSA & cases \\
ECB policy surprises (pp) & \citet{jarocinski2020information} & MP\_median (OIS) & LP-IV \\
Oil supply news (\%) & \citet{kanzig2021oil} & surprise series & LP-IV \\
Value added in final demand (USD) & \citet{oecd2025tiva} & FDVA & shift-share \\
Employment protection (index 0--6) & \citet{oecd2020epl} & EPL\_R, EPL\_T, version 3, 2008 & cases \\
\bottomrule\end{tabularx}
\end{table}

\subsection{Measurement issues}

\paragraph{Domestic and resident concepts.} GDP measures production by resident
producer units, including work by cross-border commuters. Unemployment and participation
describe residents. The two diverge most in Luxembourg, where commuters fill a large share
of domestic jobs; section~\ref{sec:cases} tests whether this explains Luxembourg's weak
association.

\paragraph{Survey changes.} The EU Labour Force Survey changed its legal framework in 2021.
Eurostat back-casts the headline series from 2009, but:
\begin{itemize}
  \item Spain and France are flagged ``definition differs'' from 2021-Q1, because they still
    test workers on temporary lay-off under the earlier rule;
  \item the Netherlands phased the new survey in during 2021 without a flag;
  \item detailed survey series carry breaks, for example at 2014-Q1 and 2021-Q1.
\end{itemize}
Eurostat's quality flags are stored with every observation. The Bayesian models mask
likelihood terms that rely on a flagged value. The classical review uses a
flag-screened sample as a sensitivity analysis.

\paragraph{Monthly unemployment.} Eurostat's monthly unemployment rate is a direct
monthly survey estimate in only eight countries. Elsewhere it is a trend, a three-month
moving average, or quarterly survey results interpolated with registered unemployment.
Each of these puts later months into a month's value. Section~\ref{sec:identification}
therefore restricts the monthly analysis to the eight direct-estimate countries.

\paragraph{The pandemic.} The eight quarters of 2020--2021 contain the largest movements
of GDP in the sample. Across countries, they supply a median
\val{cl_pandemic_share_median} of the variation in current-quarter GDP growth in the full
sample. Large output movements were also met by job-retention schemes that kept workers
employed. Any estimate that includes 2020--2021 mainly describes the pandemic, so we
report samples with and without it and make the choice explicit in each protocol.

\paragraph{Revisions and timing.} Every series is a revised vintage. Eurostat's vintage
datasets cover GDP only for the EU and euro-area aggregates, and monthly unemployment only
from 2021, so the forecasts in section~\ref{sec:prediction} are retrospective
(pseudo-out-of-sample), not real-time.
